\begin{abstract}
Slides in lectures and technical talks often appear before the speech that
discusses them and persist for tens of seconds. This makes them different from
a synchronous video stream: a system can encode a stable slide once, cache its
semantic evidence, and reuse that state across many simultaneous speech
translation (SST) decisions. We formulate this asynchronous, persistent-slide
setting and separate four questions that prior image-conditioned evaluations
conflate: whether the correct slide content matters, whether visual structure
adds information beyond strong OCR, whether evidence is available causally at
the relevant speech time, and whether once-per-slide processing improves the
computation-aware quality--latency frontier. We propose \policyname{}, a
non-blocking evidence manager that extracts auditable text, layout, formula,
chart-relation, and emphasis cues after slide changes and retrieves a compact
packet from causal speech state. Our benchmark protocol includes matched,
stale, cross-talk, OCR, structured-VLM, text-equivalent, and synchronous-vision
controls, with cold, amortized, and on-path costs reported separately.
\todo{Add results only after the persistent-evidence kill tests pass.}
\end{abstract}
